\documentclass[11pt]{article}
\usepackage[a4paper,margin=26mm]{geometry}
\usepackage{amsmath,amssymb,mathtools}
\usepackage{unicode-math}
\setmainfont{Latin Modern Roman}
\setmathfont{Latin Modern Math}
\usepackage{xurl,hyperref}
\hypersetup{hidelinks,pdftitle={Problem statement}}
\newcommand{\sref}[2]{\hyperlink{source:#1:#2}{[S#2]}}
\newcommand{\cataloguescope}[1]{#1}
\begin{document}
\section{Hyperbolic geometry — Question 10.4}\label{statement}

\subsection{Definitions and mathematical statement}

\paragraph{Shared notation.}$\mathbb N=\{1,2,\ldots\}$, and $\mathbb Z,\mathbb Q,\mathbb R,\mathbb C$ denote the integers, rationals, reals and complex numbers. Any other conventions are specified locally.


A surface lamination $\Lambda$ is a closed subset of $M$ locally homeomorphic to a product of a disk with a closed transverse set. It is essential when it has no sphere leaf or torus leaf bounding a solid torus, and its complementary regions are irreducible, have incompressible boundary, and have no compressing monogon. Such a monogon is a compressing disk with one boundary arc on a boundary leaf and the other along a transverse interval fiber. It is genuine if some complementary region is not an interval bundle. A sublamination is a closed union of leaves. A manifold is atoroidal if it has no incompressible embedded torus other than a boundary-parallel one. Write $\widetilde M$ for its universal cover, and put a tilde on a foliation or lamination for its lift. For metric spaces, a quasi-isometry has constants $a\ge1,b\ge0$ with $a^{-1}d(x,y)-b\le d(f(x),f(y))\le a d(x,y)+b$, and every target point lies within a bounded distance of its image. Leaf geometry uses the intrinsic path metric from a metric on $M$.
\paragraph{Statement recorded in the supplied source.}
If atoroidal $M$ admits an essential lamination, does it admit an essential lamination with quasigeodesically embedded lifted leaves, necessarily genuine in the setting of the source? Here quasigeodesic leaves means that intrinsic leaf distances and ambient distances satisfy uniform quasi-isometric embedding bounds. The new lamination may differ from the one initially given. \sref{statement}{2}
\subsection{Short English statement}
Does existence of an essential lamination on an atoroidal manifold imply existence of one with quasigeodesic lifted leaves?
\subsection{Sources}
\begin{description}
\item[\hypertarget{source:statement:1}{S1}] ULAM.AI and the UnsolvedMath Contributors, \emph{UnsolvedMath: A Curated Collection of Open Mathematics Problems}, catalogue entry. CC BY 4.0. \url{https://huggingface.co/datasets/ulamai/UnsolvedMath}.
\item[\hypertarget{source:statement:2}{S2}] Danny Calegari, \emph{Problems in foliations and laminations of 3-manifolds}, version 0.78 (2002), Question 10.4; Definitions 1.1--1.10 and local definitions and remarks. \url{https://arxiv.org/abs/math/0209081v1}.
\end{description}
\label{end:statement}
\end{document}
